% appendix-numerical-general.tex — \input'd from appendix-numerical.tex
%
% General approach to numerical implementation:
% three-valued predicate framework and information flow architecture.
%
% Sources: restructured from appendix-numerical-three-valued.tex (lines 9-71),
%   with new information flow section based on Jörn's design (2026-02-24).
%
% Structure:
%   - Soundness goal
%   - Three-valued predicates (scoped to β>0)
%   - Impact on control flow (inclusion decisions, maximum tracking)

The numerical algorithm must never report a capacity value
that is provably wrong: if the algorithm certifies a value,
that value must correspond to a genuine critical point of the
exact problem.
We formalise this using three-valued predicates.

\begin{definition}[Sound predicate]%
\label{def:sound-predicate}
A predicate~$\tilde P$ computed from floating-point data
is \emph{sound} if its definite verdicts are reliable:
$\textsc{true}$ and $\textsc{false}$ agree with exact
arithmetic, and $\textsc{unknown}$ means the numerical
evidence is inconclusive.
Formally:
\begin{align*}
  P \text{ true in exact arithmetic}
  &\;\Longrightarrow\;
  \tilde P \neq \textsc{false}, \\
  P \text{ false in exact arithmetic}
  &\;\Longrightarrow\;
  \tilde P \neq \textsc{true}.
\end{align*}
\end{definition}

Not every predicate in the implementation needs this
treatment---many decisions (e.g.\ index comparisons,
loop bounds) are exact in floating-point arithmetic.
The three-valued framework applies to predicates
whose evaluation involves numerical error, principally
the positivity predicate $\beta_i > 0$
(Section~\ref{app:svd-step}).

\begin{remark}[Impact on control flow]%
\label{rem:three-valued-control-flow}
The exact algorithm has two kinds of decision points:
\begin{description}
\item[Inclusion decisions] (``unless $P$, skip''):
  The exact rule ``if $\lnot P$, skip $(S, \sigma)$''
  becomes ``if $\tilde P = \textsc{false}$, skip.''
  When $\tilde P = \textsc{unknown}$, the pair is kept
  (conservatively), since the exact~$P$ might be true.
  This means the numerical algorithm searches at least as
  broadly as the exact algorithm.
\item[Maximum tracking] (``find the largest admissible $Q$''):
  An admissibility predicate $A(\beta)$ (e.g.\ $\beta_i > 0$
  for all~$i$) becomes three-valued.
  The algorithm returns:
  \begin{itemize}
  \item the largest~$Q$ where $\tilde A = \textsc{true}$
    (the \emph{certified value}, a lower bound on the
    true maximum),
  \item all~$Q$ values above the certified value where
    $\tilde A = \textsc{unknown}$
    (the true maximum may be among these).
  \end{itemize}
  Candidates with $\tilde A = \textsc{false}$ are provably
  inadmissible and discarded; uncertain candidates at or below
  the certified value are irrelevant (the certified value already
  dominates).
\end{description}
\end{remark}

